\section*{תקציר}

מחקר זה עוסק בהתפתחות מנגנוני תיאום כלים מרובים בסוכני \textenglish{LLM}. בעשור האחרון הפכו מודלי שפה למרכיב מרכזי במערכות בינה מלאכותית, ועלה הצורך לתאם ביניהם לבין כלים חיצוניים כגון מנועי חיפוש, מחשבונים ובסיסי נתונים. המחקר בוחן שלוש גישות עיקריות: שיטות היוריסטיות מבוססות-דמיון; מסגרות למידה מפוקחת כגון \textenglish{Toolformer} ו-\textenglish{AvaTaR}; וארכיטקטורות רב-סוכניות. בנוסף נבחן \textenglish{Model Context Protocol} (\textenglish{MCP}) כתקן פתוח לממשק כלים. הממצאים מצביעים כי מסגרות המשלבות לולאות חשיבה-פעולה (\textenglish{ReAct}, \textenglish{Reflexion}) עם תכנון מודע-תלויות (\textenglish{DEPS}) משפרות משמעותית את יכולת ההרחבה, הביצועים והתחזוקה של מערכות סוכנים ייצורניות.

\chapter{הקדמה והקשר היסטורי}

תחום הבינה המלאכותית חווה שינוי דרמטי בשנים האחרונות, הנובע בעיקר מהתפתחות בחקר למידה עמוקה ומודלים של שפה טבעית. הרקע ההיסטורי של תחום זה מסתמך על מחקר שנתגבש לאורך עשרות שנים, החל מתרגיל הודעה שנוצר בשנות החמישים של המאה העשרים.

בעבר, מערכות בינה מלאכותית היו מבוססות על כללים מפורשים שהתחזקו על ידי מומחים אנושיים. אולם, עם העלייה של שיטות למידה ממוכנת, באופן הדרגתי יותר ויותר מערכות החלו ללמוד דפוסים מהנתונים ישירות. המהלך הזה הכניס חידוש עמוק בדרך בה מחשבים מעבדים מידע טבעי, במיוחד בהקשר של שפות אנושיות.

עם הופעתם של מודלים כמו \textenglish{Transformer}, תחום עיבוד השפה הטבעית (\textenglish{NLP}) עבר מהפכה. מנגנון תשומת הלב (\textenglish{attention mechanism}) המשמש בגרעין הארכיטקטורה של \textenglish{Transformer} אפשר לדגמים להתמקד בחלקים רלוונטיים של הקלט בזמן עיבודו. זהו שינוי משמעותי מטכניקות קודמות כמו רשתות נוירונים רקורנטיות (\textenglish{RNNs}).

\section{בקרת הקשר ומודלים בגדול}

היום, מודלים גדולי של שפה (\textenglish{Large Language Models}, או \textenglish{LLMs}) כמו \textenglish{GPT-3}, \textenglish{Claude}, ו-\textenglish{LLaMA} שינו את הנוף של יישומים מיישומיים ועסקיים. מודלים אלה מאומנים על כמויות עצומות של נתוני טקסט, וייצורים שלהם כוללים תוכן מוסדר, תרגומים, עיבוד שאלות ותשובות, וכלים אחרים של בינה מלאכותית.

עם כל התקדמות בחוזק וביכולת של דגמים, עולות גם סוגיות חדשות של בקרה והשגחה. כיצד יכולות ארגונים לבקר על הדגמים האלה? כיצד יכולנו להבטיח שהם עובדים לפי הערכים וההנחיות של חברות וחברה בכללותה?

פרק זה מנסה להציע מסגרת מעשית למשימות אלה, עם דגש על בקרת הקשר (\textenglish{context control}) ומודל פרוטוקול של תקשורת (\textenglish{Model Context Protocol}). בעזרת כלים אלה, ארגונים יכולים להשגיח ביעילות על שיטות שלהם בעבודה עם \textenglish{LLMs}, וכך לשיפר את כמות הביקורת הזמינה בעת שימוש במערכות אלו בסביבות עסקיות ומקצועיות.
